%!TEX program = xelatex
\documentclass[12pt,a4paper]{article}
\usepackage[margin=21mm,top=25mm,bottom=23mm]{geometry}
\usepackage{fontspec,polyglossia,amsmath,amssymb,mathtools,microtype}
\usepackage{tikz,circuitikz,xcolor,hyperref,fancyhdr,tcolorbox,needspace,enumitem}
\setmainlanguage{arabic}
\setotherlanguage{english}
\newfontfamily\arabicfont[Script=Arabic,Scale=1.08]{Amiri}
\newfontfamily\arabicfontsf[Script=Arabic,Scale=1.00]{Amiri}
\setmainfont{DejaVu Serif}
\setsansfont{DejaVu Sans}
\definecolor{ink}{HTML}{213C36}
\definecolor{accent}{HTML}{187767}
\definecolor{muted}{HTML}{64776F}
\definecolor{linecol}{HTML}{E1E9E5}
\definecolor{light}{HTML}{F1F7F3}
\definecolor{amber}{HTML}{B47B3D}
\hypersetup{colorlinks=true,urlcolor=accent,linkcolor=accent,pdfauthor={Control Theory Study Notes},pdftitle={حل واجبات المحاضرة الاولى}}
\setlength{\parindent}{0pt}
\setlength{\parskip}{5pt}
\renewcommand{\baselinestretch}{1.13}
\setlist[itemize]{topsep=2pt,itemsep=2pt,parsep=0pt,rightmargin=16pt}
\pagestyle{fancy}
\fancyhf{}
\fancyhead[R]{\footnotesize\textcolor{muted}{حل واجبات المحاضرة الاولى}}
\fancyhead[L]{\footnotesize\textcolor{muted}{\textenglish{CONTROL THEORY}}}
\fancyfoot[C]{\footnotesize\textcolor{muted}{\thepage}}
\renewcommand{\headrulewidth}{0.3pt}
\renewcommand{\footrulewidth}{0pt}
\setlength{\headheight}{17pt}
\newcommand{\M}[1]{\begingroup\begin{LTR}\begin{equation*}#1\end{equation*}\end{LTR}\endgroup}
\newcommand{\Q}[1]{\begingroup\small\begin{LTR}\begin{equation*}#1\end{equation*}\end{LTR}\endgroup}
\newcommand{\Step}[2]{\Needspace{65pt}\vspace{8pt}{\arabicfontsf\bfseries\large\color{ink} #2}\par\vspace{2pt}\begin{LTR}\noindent{\color{accent}\footnotesize\sffamily STEP #1}\end{LTR}\noindent\textcolor{linecol}{\rule{\textwidth}{0.5pt}}\par}
\newcommand{\Homework}[1]{\clearpage\vspace*{5pt}\begin{LTR}\noindent\textcolor{accent}{\rule{34mm}{2.2pt}}\\[7pt]{\sffamily\bfseries\Large Homework #1}\end{LTR}\vspace{4pt}\noindent{\LARGE\arabicfontsf\bfseries\color{ink}الحل بطريقة \textenglish{Mesh Analysis}}\par\vspace{3pt}}
\newcommand{\Final}[1]{\par\vspace{9pt}\begin{tcolorbox}[colback=light,colframe=accent,boxrule=0.7pt,arc=2mm,left=12pt,right=12pt,top=11pt,bottom=9pt]\noindent{\large\bfseries\textcolor{accent}{\textenglish{Final Answer} — الجواب النهائي}}\par\M{\boxed{#1}}\end{tcolorbox}}
\newcommand{\Remember}[1]{\vspace{10pt}\noindent\textbf{للحفظ والفهم:} #1}
\AtBeginDocument{\color{ink}}
\begin{document}
\thispagestyle{empty}
\vspace*{33mm}
\begin{center}
{\color{accent}\textenglish{CONTROL THEORY}\\[5pt]\textenglish{ELECTRICAL SYSTEMS}}\\[26pt]
{\arabicfontsf\bfseries\fontsize{27}{41}\selectfont حل واجبات\\المحاضرة الاولى}\\[12pt]
{\color{muted}\arabicfont\Large طريقة تيارات الحلقات — \textenglish{Mesh Analysis}}\\[14pt]
\begin{tikzpicture}\draw[accent,line width=1.2pt](-3.5,0)--(3.5,0);\fill[amber](0,0)circle(2pt);\end{tikzpicture}\\[21pt]
{\textenglish{Homework 01 / 02 / 03}}\\[5pt]
{\large الأسئلة بالعربية والإنجليزية، والحلول الرياضية خطوة بخطوة}
\end{center}
\vfill
\begin{tcolorbox}[colback=light,colframe=linecol,boxrule=.5pt,arc=2mm]
\begin{center}
\textbf{مادة نظرية السيطرة — المرحلة الثالثة}\\
\textenglish{Transfer Functions • Kirchhoff's Laws • Laplace Impedances}
\end{center}
\end{tcolorbox}
\newpage
\Homework{1}
\textbf{السؤال بالإنجليزي (\textenglish{English}):}
\begin{LTR}\begin{english}
Find the transfer function \(\dfrac{V_o(s)}{V_i(s)}\) of the following network.
\end{english}\end{LTR}
\textbf{السؤال بالعربي:}
أوجد دالة الانتقال للدائرة الكهربائية الموضحة في الشكل باستخدام طريقة تحليل تيارات الحلقات (\textenglish{Mesh Analysis}).

المطلوب:
\M{\boxed{G(s)=\frac{V_o(s)}{V_i(s)}}}
\begin{center}
\begin{circuitikz}[american,scale=.88,transform shape]
\draw (0,0) to[sV,l=$V_i$] (0,3) to[L,l=$L_1$] (3.2,3) to[L,l=$L_2$,v^>=$V_o$] (6.4,3) to[C,l=$C$] (6.4,0)--(0,0);
\draw (3.2,3) to[R,l=$R$] (3.2,0);
\node at (1.65,1.4) {$I_1$};\node at (4.82,1.4) {$I_2$};
\end{circuitikz}
\end{center}
\Step{1}{تحديد تيارات الحلقات (\textenglish{Mesh Currents})}
نفترض تيارين باتجاه عقارب الساعة:
\begin{itemize}
\item $I_1$: تيار الحلقة الأولى، التي تحتوي المصدر $V_i$، والمحث $L_1$، والمقاومة $R$.
\item $I_2$: تيار الحلقة الثانية، التي تحتوي المقاومة $R$، والمحث $L_2$، والمتسعة $C$.
\end{itemize}
المقاومة $R$ مشتركة بين الحلقتين، ولذلك تيارها باتجاه الحلقة الأولى:
\M{\boxed{I_R=I_1-I_2}}
\Step{2}{تحويل العناصر إلى ممانعات (\textenglish{Impedances})}
\M{Z_{L_1}=L_1s\qquad Z_{L_2}=L_2s}
\M{Z_R=R\qquad Z_C=\frac{1}{Cs}}
نستخدم العلاقة:
\M{\boxed{V=ZI}}
\Step{3}{تطبيق \textenglish{KVL} على الحلقة الأولى (\textenglish{Mesh 1})}
حسب قانون كيرشوف، جهد المصدر يساوي مجموع هبوط الجهد على المحث الأول والمقاومة:
\M{V_i=V_{L_1}+V_R}
نعوض:
\M{V_i=L_1sI_1+R(I_1-I_2)}
نفتح الأقواس:
\M{V_i=L_1sI_1+RI_1-RI_2}
نجمع حدود $I_1$:
\M{\boxed{(L_1s+R)I_1-RI_2=V_i}\tag{1}}
\Step{4}{تطبيق \textenglish{KVL} على الحلقة الثانية (\textenglish{Mesh 2})}
الحلقة الثانية لا تحتوي مصدراً، لذلك:
\M{V_R+V_{L_2}+V_C=0}
نعوض:
\M{R(I_2-I_1)+L_2sI_2+\frac{I_2}{Cs}=0}
نفتح الأقواس:
\M{RI_2-RI_1+L_2sI_2+\frac{I_2}{Cs}=0}
نجمع حدود $I_2$:
\M{\boxed{-RI_1+\left(R+L_2s+\frac{1}{Cs}\right)I_2=0}\tag{2}}
الآن لدينا معادلتان بمجهولين $I_1,I_2$.
\Step{5}{حذف التيار $I_1$ (\textenglish{Elimination})}
من المعادلة (2):
\M{-RI_1+\left(R+L_2s+\frac1{Cs}\right)I_2=0}
ننقل الحد الذي يحتوي $I_1$:
\M{RI_1=\left(R+L_2s+\frac1{Cs}\right)I_2}
نقسم على $R$:
\M{\boxed{I_1=\frac{R+L_2s+\frac1{Cs}}{R}I_2}}
نعوض في المعادلة (1):
\M{(L_1s+R)I_1-RI_2=V_i}
\M{(L_1s+R)\left(\frac{R+L_2s+\frac1{Cs}}{R}I_2\right)-RI_2=V_i}
نأخذ $I_2$ عاملاً مشتركاً:
\Q{\left[\frac{(L_1s+R)(R+L_2s+\frac1{Cs})}{R}-R\right]I_2=V_i}
نوحد المقامات:
\Q{\left[\frac{(L_1s+R)(R+L_2s+\frac1{Cs})-R^2}{R}\right]I_2=V_i}
نعزل $I_2$:
\Q{I_2=\frac{RV_i}{(L_1s+R)(R+L_2s+\frac1{Cs})-R^2}}
نقسم على $V_i$:
\Q{\boxed{\frac{I_2}{V_i}=\frac{R}{(L_1s+R)(R+L_2s+\frac1{Cs})-R^2}}\tag{3}}
\Step{6}{إيجاد جهد الخرج (\textenglish{Output Voltage})}
لأن الخرج مأخوذ على المحث الثاني $L_2$، والتيار المار فيه $I_2$:
\M{V_o=Z_{L_2}I_2}
\M{\boxed{V_o=L_2sI_2}\tag{4}}
نقسم على $V_i$:
\M{\frac{V_o}{V_i}=L_2s\frac{I_2}{V_i}}
نعوض من المعادلة (3):
\Q{\frac{V_o}{V_i}=L_2s\frac{R}{(L_1s+R)(R+L_2s+\frac1{Cs})-R^2}}
إذن:
\Q{\boxed{\frac{V_o}{V_i}=\frac{RL_2s}{(L_1s+R)(R+L_2s+\frac1{Cs})-R^2}}}
\Step{7}{تبسيط دالة الانتقال (\textenglish{Simplification})}
نفتح الأقواس في المقام:
\Q{\begin{aligned}&(L_1s+R)\left(R+L_2s+\frac1{Cs}\right)-R^2\\&=L_1Rs+L_1L_2s^2+\frac{L_1}{C}\\&\quad+R^2+RL_2s+\frac{R}{Cs}-R^2\end{aligned}}
نختصر $R^2-R^2$:
\M{=L_1Rs+L_1L_2s^2+\frac{L_1}{C}+RL_2s+\frac{R}{Cs}}
نجمع الحدود المتشابهة:
\M{=L_1L_2s^2+R(L_1+L_2)s+\frac{L_1}{C}+\frac{R}{Cs}}
إذن:
\Q{\frac{V_o}{V_i}=\frac{RL_2s}{L_1L_2s^2+R(L_1+L_2)s+\frac{L_1}{C}+\frac{R}{Cs}}}
نضرب البسط والمقام في $Cs$:
\Q{\frac{V_o}{V_i}=\frac{RL_2Cs^2}{L_1L_2Cs^3+RC(L_1+L_2)s^2+L_1s+R}}
وأخيراً نقسم البسط والمقام على $R$.
\Final{\frac{V_o(s)}{V_i(s)}=\frac{L_2Cs^2}{\frac{L_1L_2C}{R}s^3+(L_1+L_2)Cs^2+\frac{L_1}{R}s+1}}
\Remember{نكتب معادلتي \textenglish{KVL} للحلقتين، ثم نحذف $I_1$ لنحصل على $I_2$. بعدها نستخدم $V_o=L_2sI_2$، لأن الخرج مقاس على المحث الثاني، ونستخرج دالة الانتقال.}
\Homework{2}
\textbf{السؤال بالإنجليزي (\textenglish{English}):}
\begin{LTR}\begin{english}
Find the transfer function \(\dfrac{E_o(s)}{E_i(s)}\) of the following network.
\end{english}\end{LTR}
\textbf{السؤال بالعربي:}
أوجد دالة الانتقال للدائرة الكهربائية الموضحة في الشكل باستخدام طريقة تحليل تيارات الحلقات (\textenglish{Mesh Analysis}).
المطلوب:
\M{\boxed{G(s)=\frac{E_o(s)}{E_i(s)}}}
تحتوي الدائرة على مقاومتين $R_1,R_2$ ومتسعتين $C_1,C_2$. المقاومة $R_2$ والمتسعة $C_2$ متصلتان على التوازي، بينما المقاومة $R_1$ والمتسعة $C_1$ متصلتان على التوالي في الفرع الأيمن.
جهد الخرج مأخوذ على المتسعة $C_1$ فقط.
\begin{center}
\begin{circuitikz}[american,scale=.84,transform shape]
\draw (0,0) to[sV,l=$E_i$] (0,3) --(1.2,3) --(1.2,4.2) to[R,l=$R_2$] (4.6,4.2) -- (4.6,3) -- (6,3) to[R,l=$R_1$] (6,1.6) to[C,l=$C_1$,v^>=$E_o$] (6,0) -- (0,0);
\draw (1.2,3) to[C,l=$C_2$] (4.6,3);
\node at (2.4,1.5) {$I_1$};\node at (3,3.72) {$I_2$};
\end{circuitikz}
\end{center}
\Step{1}{تحديد تيارات الحلقات (\textenglish{Mesh Currents})}
نفترض تيارين باتجاه عقارب الساعة:
\begin{itemize}
\item $I_1$: تيار الحلقة السفلية، التي تحتوي المصدر $E_i$، والمتسعة $C_2$، والمقاومة $R_1$، والمتسعة $C_1$.
\item $I_2$: تيار الحلقة العلوية، التي تحتوي المقاومة $R_2$ والمتسعة $C_2$.
\end{itemize}
المتسعة $C_2$ مشتركة بين الحلقتين، ولذلك تيارها باتجاه الحلقة الأولى:
\M{\boxed{I_{C_2}=I_1-I_2}}
أما المقاومة $R_1$ والمتسعة $C_1$، فيمر فيهما تيار الحلقة الأولى فقط $I_1$.
\Step{2}{تحويل العناصر إلى ممانعات (\textenglish{Impedances})}
\M{Z_{R_1}=R_1\qquad Z_{R_2}=R_2}
\M{Z_{C_1}=\frac{1}{C_1s}\qquad Z_{C_2}=\frac{1}{C_2s}}
نستخدم العلاقة:
\M{\boxed{V=ZI}}
\Step{3}{تطبيق \textenglish{KVL} على الحلقة الأولى (\textenglish{Mesh 1})}
حسب قانون كيرشوف، جهد المصدر يساوي مجموع هبوط الجهد على المتسعة $C_2$، والمقاومة $R_1$، والمتسعة $C_1$:
\M{E_i=V_{C_2}+V_{R_1}+V_{C_1}}
نعوض:
\M{E_i=\frac{I_1-I_2}{C_2s}+R_1I_1+\frac{I_1}{C_1s}}
نفتح الأقواس:
\M{E_i=\frac{I_1}{C_2s}-\frac{I_2}{C_2s}+R_1I_1+\frac{I_1}{C_1s}}
نجمع حدود $I_1$:
\Q{\boxed{\left(R_1+\frac{1}{C_1s}+\frac{1}{C_2s}\right)I_1-\frac{I_2}{C_2s}=E_i}\tag{1}}
\Step{4}{تطبيق \textenglish{KVL} على الحلقة الثانية (\textenglish{Mesh 2})}
الحلقة العلوية لا تحتوي مصدراً، لذلك:
\M{V_{R_2}+V_{C_2}=0}
نعوض:
\M{R_2I_2+\frac{I_2-I_1}{C_2s}=0}
نفتح الأقواس:
\M{R_2I_2+\frac{I_2}{C_2s}-\frac{I_1}{C_2s}=0}
نجمع حدود $I_2$:
\Q{\boxed{-\frac{I_1}{C_2s}+\left(R_2+\frac{1}{C_2s}\right)I_2=0}\tag{2}}
الآن لدينا معادلتان بمجهولين $I_1,I_2$.
\Step{5}{حذف التيار $I_2$ (\textenglish{Elimination})}
في هذا الواجب نريد في النهاية $I_1$، لأن جهد الخرج مأخوذ على المتسعة $C_1$، والتيار المار فيها هو $I_1$.
من المعادلة (2):
\M{-\frac{I_1}{C_2s}+\left(R_2+\frac{1}{C_2s}\right)I_2=0}
نضرب الطرفين في $C_2s$:
\M{-I_1+(R_2C_2s+1)I_2=0}
ننقل $I_1$ للطرف الآخر:
\M{I_1=(R_2C_2s+1)I_2}
نقسم على $R_2C_2s+1$:
\M{\boxed{I_2=\frac{I_1}{R_2C_2s+1}}\tag{3}}
نعوض في المعادلة (1):
\Q{\left(R_1+\frac{1}{C_1s}+\frac{1}{C_2s}\right)I_1-\frac{I_2}{C_2s}=E_i}
\Q{\left(R_1+\frac{1}{C_1s}+\frac{1}{C_2s}\right)I_1-\frac{1}{C_2s}\left(\frac{I_1}{R_2C_2s+1}\right)=E_i}
نأخذ $I_1$ عاملاً مشتركاً:
\Q{\left[R_1+\frac{1}{C_1s}+\frac{1}{C_2s}-\frac{1}{C_2s(R_2C_2s+1)}\right]I_1=E_i}
نجمع الحدين المتعلقين بالمتسعة $C_2$:
\M{\frac{1}{C_2s}-\frac{1}{C_2s(R_2C_2s+1)}}
نوحد المقامين:
\M{=\frac{R_2C_2s+1-1}{C_2s(R_2C_2s+1)}}
\M{=\frac{R_2C_2s}{C_2s(R_2C_2s+1)}}
نختصر $C_2s$:
\M{=\frac{R_2}{R_2C_2s+1}}
إذن معادلة الدخل تصبح:
\Q{\left[R_1+\frac{1}{C_1s}+\frac{R_2}{R_2C_2s+1}\right]I_1=E_i}
نعزل $I_1$:
\M{I_1=\frac{E_i}{R_1+\frac{1}{C_1s}+\frac{R_2}{R_2C_2s+1}}}
نقسم على $E_i$:
\Q{\boxed{\frac{I_1}{E_i}=\frac{1}{R_1+\frac{1}{C_1s}+\frac{R_2}{R_2C_2s+1}}}\tag{4}}
\Step{6}{إيجاد جهد الخرج (\textenglish{Output Voltage})}
لأن الخرج مأخوذ على المتسعة $C_1$، والتيار المار فيها $I_1$:
\M{E_o=V_{C_1}}
نستخدم قانون:
\M{V=ZI}
إذن:
\M{E_o=Z_{C_1}I_1}
نعوض عن ممانعة المتسعة:
\M{\boxed{E_o=\frac{I_1}{C_1s}}\tag{5}}
نقسم على $E_i$:
\M{\frac{E_o}{E_i}=\frac{1}{C_1s}\frac{I_1}{E_i}}
نعوض من المعادلة (4):
\Q{\frac{E_o}{E_i}=\frac{1}{C_1s}\times\frac{1}{R_1+\frac{1}{C_1s}+\frac{R_2}{R_2C_2s+1}}}
إذن:
\Q{\boxed{\frac{E_o}{E_i}=\frac{\frac{1}{C_1s}}{R_1+\frac{1}{C_1s}+\frac{R_2}{R_2C_2s+1}}}}
\Step{7}{تبسيط دالة الانتقال (\textenglish{Simplification})}
نضرب البسط والمقام في $C_1s$:
\M{\frac{E_o}{E_i}=\frac{1}{R_1C_1s+1+\frac{R_2C_1s}{R_2C_2s+1}}}
لإزالة الكسر الموجود في المقام، نضرب البسط والمقام في $R_2C_2s+1$:
\Q{\frac{E_o}{E_i}=\frac{R_2C_2s+1}{(R_1C_1s+1)(R_2C_2s+1)+R_2C_1s}}
نفتح الأقواس:
\Q{\begin{aligned}&(R_1C_1s+1)(R_2C_2s+1)\\&=R_1R_2C_1C_2s^2+R_1C_1s+R_2C_2s+1\end{aligned}}
نضيف الحد الموجود خارج الأقواس:
\Q{\begin{aligned}&=R_1R_2C_1C_2s^2+R_1C_1s\\&\quad+R_2C_2s+R_2C_1s+1\end{aligned}}
نجمع حدود $s$:
\Q{=R_1R_2C_1C_2s^2+(R_1C_1+R_2C_2+R_2C_1)s+1}
\Final{\frac{E_o(s)}{E_i(s)}=\frac{R_2C_2s+1}{R_1R_2C_1C_2s^2+(R_1C_1+R_2C_2+R_2C_1)s+1}}
\Remember{نكتب معادلتي \textenglish{KVL}، ثم نحذف $I_2$ للحصول على $I_1$، لأن التيار في المتسعة $C_1$ هو $I_1$. بعدها نستخدم $E_o=I_1/(C_1s)$ ونقسم على $E_i$ لاستخراج دالة الانتقال.}
\Homework{3}
\textbf{السؤال بالإنجليزي (\textenglish{English}):}
\begin{LTR}\begin{english}
Find the transfer function \(\dfrac{I_2(s)}{E_i(s)}\) of the following network.
\end{english}\end{LTR}
\textbf{السؤال بالعربي:}
أوجد دالة الانتقال للدائرة الكهربائية الموضحة في الشكل باستخدام طريقة تحليل تيارات الحلقات (\textenglish{Mesh Analysis}).
المطلوب:
\M{\boxed{G(s)=\frac{I_2(s)}{E_i(s)}}}
تحتوي الدائرة على مقاومتين $R_1,R_2$، ومحث $L$، ومتسعة $C$.
لاحظ أن المطلوب هذه المرة هو نسبة تيار الحلقة الثانية $I_2$ إلى جهد الدخل $E_i$، وليس جهد الخرج إلى جهد الدخل.
\begin{center}
\begin{circuitikz}[american,scale=.87,transform shape]
\draw (0,0) to[sV,l=$E_i$] (0,3) to[R,l=$R_1$] (3.2,3) to[R,l=$R_2$] (6.4,3) to[C,l=$C$] (6.4,0)--(0,0);
\draw (3.2,3) to[L,l=$L$] (3.2,0);
\node at (1.6,1.35) {$I_1$};\node at (4.85,1.35) {$I_2$};
\end{circuitikz}
\end{center}
\Step{1}{تحديد تيارات الحلقات (\textenglish{Mesh Currents})}
نفترض تيارين باتجاه عقارب الساعة:
\begin{itemize}
\item $I_1$: تيار الحلقة الأولى، التي تحتوي المصدر $E_i$، والمقاومة $R_1$، والمحث $L$.
\item $I_2$: تيار الحلقة الثانية، التي تحتوي المحث $L$، والمقاومة $R_2$، والمتسعة $C$.
\end{itemize}
المحث $L$ مشترك بين الحلقتين، ولذلك تياره باتجاه الحلقة الأولى:
\M{\boxed{I_L=I_1-I_2}}
أما باتجاه الحلقة الثانية:
\M{\boxed{I_L=I_2-I_1}}
\Step{2}{تحويل العناصر إلى ممانعات (\textenglish{Impedances})}
\M{Z_{R_1}=R_1\qquad Z_{R_2}=R_2}
\M{Z_L=Ls\qquad Z_C=\frac{1}{Cs}}
نستخدم العلاقة:
\M{\boxed{V=ZI}}
\Step{3}{تطبيق \textenglish{KVL} على الحلقة الأولى (\textenglish{Mesh 1})}
حسب قانون كيرشوف، جهد المصدر يساوي مجموع هبوط الجهد على المقاومة الأولى والمحث:
\M{E_i=V_{R_1}+V_L}
نعوض:
\M{E_i=R_1I_1+Ls(I_1-I_2)}
نفتح الأقواس:
\M{E_i=R_1I_1+LsI_1-LsI_2}
نجمع حدود $I_1$:
\M{\boxed{(R_1+Ls)I_1-LsI_2=E_i}\tag{1}}
\Step{4}{تطبيق \textenglish{KVL} على الحلقة الثانية (\textenglish{Mesh 2})}
الحلقة الثانية لا تحتوي مصدراً، لذلك:
\M{V_L+V_{R_2}+V_C=0}
نعوض:
\M{Ls(I_2-I_1)+R_2I_2+\frac{I_2}{Cs}=0}
نفتح الأقواس:
\M{LsI_2-LsI_1+R_2I_2+\frac{I_2}{Cs}=0}
نرتب الحدود:
\M{-LsI_1+LsI_2+R_2I_2+\frac{I_2}{Cs}=0}
نجمع حدود $I_2$:
\Q{\boxed{-LsI_1+\left(Ls+R_2+\frac{1}{Cs}\right)I_2=0}\tag{2}}
الآن لدينا معادلتان بمجهولين $I_1,I_2$.
\Step{5}{حذف التيار $I_1$ (\textenglish{Elimination})}
نريد استخراج $I_2$، لذلك نحذف $I_1$.
من المعادلة (2):
\M{-LsI_1+\left(Ls+R_2+\frac1{Cs}\right)I_2=0}
ننقل الحد الذي يحتوي $I_1$ إلى الطرف الآخر:
\M{\left(Ls+R_2+\frac1{Cs}\right)I_2=LsI_1}
نقسم على $Ls$:
\M{\boxed{I_1=\frac{Ls+R_2+\frac1{Cs}}{Ls}I_2}\tag{3}}
نعوض في المعادلة (1):
\M{(R_1+Ls)I_1-LsI_2=E_i}
\Q{(R_1+Ls)\left(\frac{Ls+R_2+\frac1{Cs}}{Ls}I_2\right)-LsI_2=E_i}
نأخذ $I_2$ عاملاً مشتركاً:
\Q{\left[\frac{(R_1+Ls)(Ls+R_2+\frac1{Cs})}{Ls}-Ls\right]I_2=E_i}
نوحد المقامات:
\M{Ls=\frac{L^2s^2}{Ls}}
إذن:
\Q{\left[\frac{(R_1+Ls)(Ls+R_2+\frac1{Cs})-L^2s^2}{Ls}\right]I_2=E_i}
نضرب الطرفين في $Ls$:
\Q{\left[(R_1+Ls)(Ls+R_2+\frac1{Cs})-L^2s^2\right]I_2=LsE_i}
نقسم على القوس:
\Q{I_2=\frac{LsE_i}{(R_1+Ls)(Ls+R_2+\frac1{Cs})-L^2s^2}}
نقسم على $E_i$:
\Q{\boxed{\frac{I_2}{E_i}=\frac{Ls}{(R_1+Ls)(Ls+R_2+\frac1{Cs})-L^2s^2}}\tag{4}}
هذه هي دالة الانتقال قبل التبسيط.
\Step{6}{تبسيط دالة الانتقال (\textenglish{Simplification})}
نبدأ بفتح الأقواس في المقام:
\Q{(R_1+Ls)\left(Ls+R_2+\frac1{Cs}\right)-L^2s^2}
نضرب كل حد من القوس الأول في حدود القوس الثاني:
\Q{\begin{aligned}={}&R_1(Ls)+R_1(R_2)+R_1\left(\frac1{Cs}\right)\\&+Ls(Ls)+Ls(R_2)+Ls\left(\frac1{Cs}\right)-L^2s^2\end{aligned}}
ننفذ عمليات الضرب:
\Q{\begin{aligned}={}&R_1Ls+R_1R_2+\frac{R_1}{Cs}\\&+L^2s^2+LR_2s+\frac{L}{C}-L^2s^2\end{aligned}}
نختصر:
\M{+L^2s^2-L^2s^2=0}
فتبقى:
\M{=R_1Ls+R_1R_2+\frac{R_1}{Cs}+LR_2s+\frac{L}{C}}
نجمع حدود $s$:
\M{R_1Ls+LR_2s=L(R_1+R_2)s}
إذن يصبح المقام:
\M{=L(R_1+R_2)s+R_1R_2+\frac{R_1}{Cs}+\frac{L}{C}}
نعوض في دالة الانتقال:
\Q{\frac{I_2}{E_i}=\frac{Ls}{L(R_1+R_2)s+R_1R_2+\frac{R_1}{Cs}+\frac{L}{C}}}
\Step{7}{إزالة الكسور من المقام}
نضرب البسط والمقام في $Cs$:
\Q{\frac{I_2}{E_i}=\frac{Ls(Cs)}{\left[L(R_1+R_2)s+R_1R_2+\frac{R_1}{Cs}+\frac{L}{C}\right]Cs}}
نضرب البسط:
\M{Ls(Cs)=LCs^2}
نضرب كل حد في المقام:
\M{L(R_1+R_2)s(Cs)=LC(R_1+R_2)s^2}
\M{R_1R_2(Cs)=R_1R_2Cs}
\M{\frac{R_1}{Cs}(Cs)=R_1}
\M{\frac{L}{C}(Cs)=Ls}
إذن:
\Q{\frac{I_2}{E_i}=\frac{LCs^2}{LC(R_1+R_2)s^2+R_1R_2Cs+R_1+Ls}}
نجمع الحدود التي تحتوي $s$:
\M{R_1R_2Cs+Ls=(R_1R_2C+L)s}
ونرتب المقام حسب قوى $s$:
\Q{\frac{I_2}{E_i}=\frac{LCs^2}{LC(R_1+R_2)s^2+(R_1R_2C+L)s+R_1}}
\Final{\frac{I_2(s)}{E_i(s)}=\frac{LCs^2}{(R_1+R_2)LCs^2+(R_1R_2C+L)s+R_1}}
\Remember{نكتب معادلتي \textenglish{KVL} للحلقتين، ثم نحذف $I_1$ حتى نحصل على $I_2$ بدلالة $E_i$. في هذا الواجب لا نحتاج إلى حساب جهد الخرج $E_o$، لأن المطلوب هو $I_2/E_i$ مباشرة. لذلك نستخرج النسبة ونبسطها للحصول على دالة الانتقال النهائية.}
\end{document}